Como ya se ha comentado previamente, \textit{RTXI} es una aplicación para realizar experimentos en tiempo real usada por investigadores de la Escuela Politécnica Superior de la Universidad Autónoma de Madrid. Por ello se ha trabajado sobre la misma, desarrollando dos módulos para que puedan ser usados con el kernel de Linux \textit{Preempt-RT}. \textit{RTXI} ya cuenta con una version para \textit{POSIX} (aunque esta no puede operar en restricciones temporales muy estrictas), por lo que la implementación ha sido relativamente sencilla, implementando un módulo para la gestión de hilos \textit{realtime} `rtos\_preempt\_rt.cpp`, siguiendo los pasos indicados en la documentación de la Linux Foundation \cite{TheLinuxFundation2023}. El otro módulo, `fifo\_preempt\_rt.cpp`, no difiere mucho de la implementación de \textit{POSIX} en el proyecto (`fifo\_posix.cpp`), pero se ha dejado creado por si en un futuro se necesita especializar algún parámetro en las colas de tiempo real.

Además de esto, se ha modificado el script de instalación y los archivos CMakeLists.txt para que se pueda seleccionar la nueva versión de \textit{RTXI} (Repositorio \ref{SEC:RTXIPREEMPTREPO}). 

No obstante, ha habido algunas dificultades para poder hacer posible el uso de dicha herramienta en su versión más reciente (3.2.X), pues el programa ha sufrido algunos cambios críticos que se comentarán a continuación.
